\section{Literature Review}
\addtocontents{toc}{\protect\setcounter{tocdepth}{1}}

\subsection{Anomaly Detection in Telecommunications Networks}

Telecommunication anomaly detection is difficult because labels are scarce.
Moreover, the industry produces large amounts of data, with few viable labels that can be identified as anomalies.
Often, the problem is not detecting the anomaly but selecting which subsection of the network window to monitor.

\textcite{gonzalez2024one} apply a Dilated Convolutional Variational Autoencoder (DC-\ac{VAE}), a deep model that scores anomalies jointly across several related time series, to mobile-network counters.
They evaluate it on a fully aggregated view where labels are sparse: of 730{,}000 samples, only 9{,}000 were labeled, using a fixed set of 12 features.
Range-based F1 (range-F1) measures overlap between predicted and labeled anomaly intervals rather than isolated timestamps; they report a mean range-F1 of around 52\%, which is only modest.

\textcite{burgueno2020online} indicate that the network they analyzed had 25{,}000 mobile base station coverage areas.
Moreover, for each cell they analyzed 240 hourly \ac{KPI}s, but the study focuses on only 15 of the most important.
Detection is an online threshold state machine (normal / border / anomalous) run at cell level on that manually selected subset; they report capturing almost 80\% of the labeled network anomalies.
From this perspective, manual selection is a bottleneck because the 225 unmonitored features may still contain relevant information; for example, anomalies may occur in an unmonitored signal.

Call Detail Records (\ac{CDR}) and mobile activity records are recognized for analysis in telecom monitoring.
\textcite{sultan2018cdr} report no detection metrics.

\textcite{lemaire2024early} study early detection of urban events from Orange France mobile Network Signalling Data in Paris: cell-level, one-minute counts of network services (calls, SMS, and related signalling).
They use a Signature detector and hand-designed categorical aggregations over these service types, not temporal aggregation.
Against a curated set of documented urban events, their best precision is 28.63\%, with event recall of about 36\% at the corresponding operating point, indicating many false alarms even when the signal is useful.

Existing work assumes a fixed aggregation and therefore may miss anomalies that would only appear under a different grouping.
Automated search is useful when the best aggregation is not obvious in advance; in those papers, there is a fixed review of aggregation, but they do not address how to automate the aggregation of time series.
The next section will be about this topic.

\subsection{Anomaly Detection in Grouped Time Series}

Anomaly detection in grouped time series treats the detector input as a chosen representation, not the whole network.
The dataset is aggregated by related features/entities at the aggregation level, rather than treating a network as a flat collection of features.
It is important because the detection quality is dependent on the aggregation of features, not only the detector.
Across the reviewed studies, datasets, labels, and metrics differ, so reported scores are not directly comparable.

\textcite{mefleh2025beyond} used an unsupervised Mahalanobis distance approach in a telecommunications setting.
They report \ac{AUC} rising from 0.77 on the raw \ac{KPI} features to 0.86 when denominator information is included (for example, attempt counts that normalise ratios), then to 0.93 with adjusted ratios, and finally to 0.99 with learned adjustments.
This result shows that detection quality changes with feature transformation and aggregation.

\textcite{hua2022genad} show that China Mobile wireless base stations share some operational behavior, which GenAD, a generative reconstruction model, can learn and then use to detect anomalies on other stations in the same network.
GenAD was pre-trained on 3{,}000 unlabeled base stations, using multivariate time series features.
Their F1 score is 0.953 (precision 0.910).
That F1 uses point-adjusted segment scoring: if any timestamp inside a labeled anomaly interval is flagged, the whole interval counts as detected, so the score is not a timestamp-wise F1.

In the study by \textcite{oguz2022anomaly}, they simulate a multi-tier network that includes aggregations from the user equipment measurements and classifies the state of each femtocell access point.
Performance changes with the region of interest (\ac{ROI}) radius, signal shadowing, and aggregation methods.
At a 30\,m \ac{ROI} with 2\,dB shadowing, the 1D \ac{CNN} achieved 80.1\% accuracy, but as \ac{ROI} size and the shadowing increase, the accuracy decreases, so this shows that including more observations is not necessarily beneficial when added signals are less relevant or noisier.

Based on the reviewed studies, we can extract that grouping can strongly affect anomaly detection performance.
Although aggregations are selected manually before modeling the relevant \ac{KPI}s, entities, aggregation levels, and spatial boundaries, none of them systematically search for an aggregation structure that best characterizes the anomaly.

\subsection{Anomaly Detection in Hierarchical Time Series}

Hierarchy in anomaly detection is important because an alert such as ``traffic is abnormal'' is not enough for a telecom operator to act on.
The operator must determine whether it affects the entire network, a region or specific cell of the switch, or only one \ac{KPI} category.
So this top-down approach helps break down into smaller branches; although that hierarchy is not necessarily the best aggregation for detecting anomalies.

\textcite{hong2012tiresias} detect and localize anomalies in hierarchical \ac{ISP} call-detail data.
It analyzes customer-care call logs and signaling call records across predefined hierarchies, allowing operators to identify not only that an abnormal behavior occurs but also the level at which it appears.
They found over 94\% of reference anomalies, and separately reported 99.7\% anomaly identification accuracy.
These results were obtained from a fixed additive hierarchy and Holt--Winters spike forecasting.
So the author demonstrates hierarchical localization rather than searching over alternative aggregation structures.
It does not search for alternative partial aggregations or score non-spike discords, which is the gap addressed by this thesis.

On the other hand, \textcite{wolff2022composite} treat anomaly localization as the sequential search problem of a known tree of processes.
Instead of measuring every leaf process equally, their hierarchical dynamic search method first samples aggregate internal nodes and then moves towards the branch most likely to contain an anomaly.
This shows that hierarchy can be an active search space, not merely a reporting structure used after detection.
However, Wolff assumed tree topology is already known, as search is only within it.
This differs from our approach, which uses a \ac{GE} to generate candidate aggregations before applying the anomaly scorer.

\textcite{huang2026hydra} propose HYDRA, a multi-level hierarchy-driven anomaly detector that builds several hierarchy-based views of each series and combines detection signals across those levels.
It ranked first among 40 algorithms across 40 TSB-AD datasets.

\subsection{Matrix Profile and Time-Series Discord Detection}

Chapter~2 defined the Matrix Profile and discords.
What belongs in this chapter is the published evidence that the scorer is fast enough for many candidate series, and whether it has been used on telecommunications data.

Computational efficiency is substantially improved: a refinery time-series self-join is reduced from an estimate of 153.8 days by brute force to 1.2 hours with STAMP/STOMP \parencite{yeh2017timeseriesjoins}.
Also, \textcite{zhu2016matrixprofileii} showed GPU-STOMP could process 100 million subsequences in about 12 days, compared to the estimate of 1585 years for brute force.
STUMPY \parencite{law2019stumpy} gives the practical implementation of this process in Python.
\textcite{vanonsem2023motif} apply Matrix Profile to production \ac{KPI} series from telecom-provider customers of Skyline Communications (25 series), using it for variable-length \emph{motif} discovery rather than discord scoring.
They report the method is about 3--8 times faster than VALMOD; an example series yields both a 288-point (about one day) motif and a 12-point (about one hour) motif in a single distance-matrix pass.
No detection precision, recall, F1, or \ac{AUC} is reported (manual motif checks).
The series to analyse are chosen in advance, so the paper does not search aggregations.

\subsection{Combinatorial Optimisation for Selecting Time Series}

Little published work searches hierarchical aggregation structures.
Feature-selection methods are reviewed here instead because, from an optimisation viewpoint, the problems are similar: both search a combinatorial space of candidates and score each one with a fitness function.
The spaces themselves differ: feature selection chooses a subset of columns, whereas this dissertation searches grammar-constrained hierarchies.

For feature selection, $n$ features give $2^n$ candidate subsets, so exhaustive search is infeasible for large $n$; \textcite{xue2016survey} survey the evolutionary algorithms used for this search, including genetic algorithms, genetic programming, particle swarm optimisation, ant colony optimisation and differential evolution.

\textcite{linares2024evolutionary} use a genetic algorithm (SEISEN) whose fitness is a filter score, the Conditional Dependence Coefficient, to select lagged features for a multilayer perceptron forecaster.
On an electricity dataset, the feature set falls from 7{,}704 to 3{,}813 while \ac{MAE} improves from 207.193 to 194.400, showing that pruning the input representation can improve accuracy.

\textcite{pietron2025adnev} apply a genetic algorithm at three levels to multivariate time-series anomaly detection: it evolves subsets of input series, an autoencoder architecture for each subset, and the network weights, using validation loss as fitness.
On four sensor benchmarks (SWaT, WADI, MSL and SMAP), the resulting ensemble achieves the best point-wise F1 of 14 methods, with a mean of 0.70 against 0.54 for the next best, at a cost of up to 132 GPU-hours per dataset.

\textcite{he2021mining} search subspaces of high-dimensional financial transaction series with a greedy heuristic and an evolutionary $(\mu+\lambda)$ algorithm, scoring each candidate by discord mining.
On real data, greedy search achieves the higher mean average precision (0.86 against 0.72) but is about 15 times slower (1{,}345\,s against 88\,s).

Across these studies, the common pattern is an evolutionary search over a large candidate space with a fitness function evaluated per candidate.
None searches hierarchical aggregation structures, and none combines the search with a Matrix Profile scorer on telecommunications data; this dissertation does both with grammatical evolution.
\addtocontents{toc}{\protect\setcounter{tocdepth}{1}}
